
\setcounter{footnote}{10}

\SGARefTarget{sga3:ix:section:4}\section*{4. The Density Theorem}
\addcontentsline{toc}{section}{4. The Density Theorem}

This theorem (cf. \SGARefLink{sga3:ix:theorem:4:7}{4.7 below})%
\begingroup
\renewcommand{\thefootnote}{\(\ast\)}%
\footnote{All the results from
\SGARefLink{sga3:ix:definition:4:1}{4.1} through
\SGARefLink{sga3:ix:corollary:4:6}{4.6} are contained in EGA
IV\(_3\), 11.10, to which we refer the reader for a systematic account of
the notion of schematic density.}%
\addtocounter{footnote}{-1}%
\endgroup,
together with the theorem\footnote{Editor's note: the ``algebraicity''
theorem \SGARefLink{sga3:ix:theorem:7:1}{7.1}.} of \SGARefLink{sga3:ix:section:7}{Section~7}, will be the
essential tool in the present expos\'e and the two following ones for
passing from the infinitesimal properties of groups of multiplicative type
developed above to their ``finite'' properties.

\medskip
\noindent\SGARefTarget{sga3:ix:definition:4:1}\textbf{Definition 4.1.}
\textit{Let \(X\) be a prescheme. A family
\((Z_i)_{i\in I}\) of subpreschemes of \(X\) is said to be
\emph{schematically dense} if, for every open subset \(U\) of \(X\) and
every closed subprescheme \(U_0\) of \(U\) containing all the
\(Z_i\cap U\), one has \(U_0=U\).}

\textit{A subprescheme \(Z\) of \(X\) is said to be
\emph{schematically dense in \(X\)} if the one-member family consisting
of \(Z\) is schematically dense.}

It is immediate (cf. EGA IV\(_3\), 11.10.1) that the definition is
equivalent to saying that, for every open subset \(U\) of \(X\), every
section \(f\) of \(\mathscr O_U\) that vanishes on each \(Z_i\cap U\)
is zero. Equivalently, the intersection of the kernels of the canonical
homomorphisms
\SGARefTarget{sga3:ix:definition:4:1:display:01}
\[
  u_i:\mathscr O_X\longrightarrow (v_i)_*(\mathscr O_{Z_i})
\]
is zero, where \(v_i:Z_i\to X\) is the canonical immersion.%
\footnote{Editor's note: all the results on schematic density stated in
EGA IV\(_3\), \S11.10 follow from the results on ``separating families of
homomorphisms'' proved in \emph{loc.\ cit.}, \S11.9, to which some of the
editor's notes below will refer.}
When \(X\) lies over a prescheme \(S\), this is also equivalent to the
following assertion: for every open subset \(U\) of \(X\), every
\(S\)-prescheme \(Y\) separated over \(S\), and every pair of morphisms
\(u,v:U\to Y\) that coincide on the \(Z_i\cap U\), one has \(u=v\).
Indeed, \(u=v\) is equivalent to \(U_0=U\), where \(U_0\) is the inverse
image of the diagonal in \(Y\times_S Y\) under the \(S\)-morphism
\(U\to Y\times_S Y\) defined by \((u,v)\). The diagonal is a closed
subprescheme of \(Y\times_S Y\), so \(U_0\) is a closed subprescheme of
\(U\) containing all the \(Z_i\cap U\). Thus schematic density gives
\(U_0=U\), hence \(u=v\). Conversely, take
\(Y=\operatorname{Spec}\mathscr O_S[T]\).

In the terminology introduced at the end of EGA I, 9.5, saying that the
subprescheme \(Z\) is schematically dense in \(X\) also means that \(X\)
is the \emph{closure subprescheme} of \(Z\) in \(X\).

\medskip
\noindent\SGARefTarget{sga3:ix:lemma:4:2}\textbf{Lemma 4.2.}
\textit{Let \(X\) be a flat \(S\)-prescheme, and let
\((Z_i)_{i\in I}\) be a family of subpreschemes of \(X\), flat over
\(S\). Let \(S_0\) be a subprescheme of \(S\) defined by a nilpotent
ideal \(\mathscr J\), and suppose that the modules
\(\mathscr J^n/\mathscr J^{n+1}\) are locally free over \(S_0\).
Let \(X_0,Z_{i0}\) be obtained from \(X,Z_i\) by the base change
\(S_0\to S\). If \((Z_{i0})_{i\in I}\) is schematically dense in
\(X_0\), then \((Z_i)_{i\in I}\) is schematically dense in \(X\).}

\emph{Proof.}
Suppose \(\mathscr J^{n+1}=0\), where \(n\geq0\), and argue by induction
on \(n\); the assertion is trivial for \(n=0\). Define
\(S_m,X_m,Z_{im}\) by reduction modulo \(\mathscr J^{m+1}\), as usual.
The induction hypothesis already implies that
\((Z_{i,n-1})_{i\in I}\) is schematically dense in \(X_{n-1}\).
After replacing \(X\) by an induced open subset, it is enough to prove
that every section \(f\) of \(\mathscr O_X\) that vanishes on all the
\(Z_i\) is zero.

The section \(f_{n-1}\) of
\[
  \mathscr O_{X_{n-1}}
  =\mathscr O_X/\mathscr J^n\mathscr O_X
\]
defined by \(f\) vanishes on the \(Z_{i,n-1}\), hence is zero.
Consequently, \(f\) is a section of
\(\mathscr J^n\mathscr O_X\). Since \(X\) is flat over \(S\), there is
an isomorphism
\SGARefTarget{sga3:ix:lemma:4:2:display:01}
\[
  \mathscr J^n\mathscr O_X
  \xleftarrow{\ \sim\ }
  \mathscr E\otimes_{\mathscr O_{S_0}}\mathscr O_{X_0},
  \qquad
  \mathscr E=\mathscr J^n=\mathscr J^n/\mathscr J^{n+1}.
\]
Similarly, since each \(Z_i\) is flat over \(S\), the restriction
\(f_i\) of \(f\) to \(Z_i\) may be regarded as a section of
\SGARefTarget{sga3:ix:lemma:4:2:display:02}
\[
  \mathscr J^n\mathscr O_{Z_i}
  \xleftarrow{\ \sim\ }
  \mathscr E\otimes_{\mathscr O_{S_0}}\mathscr O_{Z_{i0}}.
\]
By hypothesis, all the \(f_i\) are zero. The module \(\mathscr E\) is
locally free by assumption, and therefore so is
\[
  \mathscr F
  =\mathscr E\otimes_{\mathscr O_{S_0}}\mathscr O_{X_0}.
\]
Thus \(f\) is a section of the locally free module \(\mathscr F\) on
\(X_0\) whose restriction to every \(Z_{i0}\) is zero. Since the family
\((Z_{i0})_{i\in I}\) is schematically dense in \(X_0\), it follows at
once that \(f=0\). \qed

\medskip
\noindent\SGARefTarget{sga3:ix:corollary:4:3}\textbf{Corollary 4.3.}
\textit{Let \(X\) be a locally Noetherian \(S\)-prescheme, flat over
\(S\), and let \((Z_i)_{i\in I}\) be a family of subpreschemes of \(X\),
flat over \(S\). Suppose that, for every \(s\in S\), the family of fibers
\((Z_{is})_{i\in I}\) is schematically dense in \(X_s\). Then
\((Z_i)_{i\in I}\) is schematically dense in \(X\).}

\emph{Proof.}
After replacing \(X\) by an induced open subset, it is enough to prove
that every section \(f\) of \(\mathscr O_X\) whose restriction to all
the \(Z_i\) is zero is itself zero.%
\footnote{Editor's note: the original has been expanded in what follows.}
For this, it suffices to show that the image of \(f\) in
\(\mathscr O_{X,x}\) is zero for every \(x\in X\). Let
\(\mathfrak m_x\) be the maximal ideal of \(\mathscr O_{X,x}\), and let
\(\mathfrak m_s\) be that of \(\mathscr O_{S,s}\), where \(s\) is the
image of \(x\) in \(S\). Since \(\mathscr O_{X,x}\) is Noetherian,
\[
  \bigcap_{n\in\mathbb N}\mathfrak m_x^n=0,
  \qquad\text{and hence, a fortiori,}\qquad
  \bigcap_{n\in\mathbb N}\mathscr O_{X,x}\mathfrak m_s^n=0.
\]
It is therefore enough to show, for every integer \(n\), that the section
induced by \(f\) on
\[
  X\times_S
  \operatorname{Spec}\bigl(\mathscr O_{S,s}/\mathfrak m_s^{n+1}\bigr)
\]
is zero. By hypothesis,
\[
  Z_{is}=Z_i\otimes_{\mathscr O_{S,s}}\kappa(s)
\]
is a schematically dense family in
\[
  X_s=X\otimes_{\mathscr O_{S,s}}\kappa(s).
\]
We are thus reduced to the case in which \(S\) is the spectrum of a
local ring whose maximal ideal \(\mathscr J\) is nilpotent,
\(S_0=\operatorname{Spec}\kappa(s)\), and the hypotheses of
\SGARefLink{sga3:ix:lemma:4:2}{Lemma~4.2} hold. The conclusion follows.
\qed

\medskip
\noindent\SGARefTarget{sga3:ix:lemma:4:4}\textbf{Lemma 4.4.}
\textit{Let \(S\) be a locally Noetherian prescheme, let \(X\) be a
locally Noetherian \(S\)-prescheme, flat over \(S\), and let
\((Z_i)_{i\in I}\) be a family of subpreschemes of \(X\), flat over
\(S\). Suppose that, for every \(s\in S\), the family of fibers
\((Z_{is})_{i\in I}\) is schematically dense in \(X_s\).
Then, for every base change \(S'\to S\), with \(S'\) not necessarily
locally Noetherian, the family \((Z_i')_{i\in I}\) is schematically
dense in \(X'\). In other words, \((Z_i)_{i\in I}\) is universally
schematically dense in \(X\) relative to \(S\) (cf. EGA IV\(_3\),
Definition 11.10.8).}

\emph{Proof.}%
\footnote{Editor's note: for another proof, reducing to the case in which
\(S'\) is locally Noetherian and using
\SGARefLink{sga3:ix:corollary:4:3}{4.3} and
\SGARefLink{sga3:ix:lemma:4:5}{4.5} as here, see EGA
IV\(_3\), 11.9.16 and 11.9.12. In the last line of the proof of 11.9.16,
replace 11.9.5 by 11.9.12.}
Let \(f'\) be a section of \(\mathscr O_{X'}\) over an open subset
\(U'\) of \(X'\), vanishing on the \(Z_i'\); we must show that \(f'=0\).
Fix \(s'\in S'\). It is enough to prove that \(f'\) vanishes at every
point of \(U'\) above \(s'\). If \(s\) is the image of \(s'\) in \(S\),
then, after replacing \(S,S'\) by the spectra of their local rings at
\(s,s'\), we may suppose that \(S,S'\) are local and that \(s,s'\) are
their closed points. We may also suppose that \(X\) is affine, so that
\SGARefTarget{sga3:ix:lemma:4:4:display:01}
\[
  \begin{gathered}
    S=\operatorname{Spec}(A),\qquad
    S'=\operatorname{Spec}(A'),\\
    X=\operatorname{Spec}(B),\qquad
    X'=\operatorname{Spec}(B'),\qquad
    B'=B\otimes_A A',
  \end{gathered}
\]
where \(A,A'\) are local, \(A\to A'\) is a local homomorphism, and
\(A,B\) are Noetherian. We may further suppose that \(U'\) is affine,
of the form \(X'_{g'}\), with \(g'\in B'\).

For every \(A\)-subalgebra \(A''\) of \(A'\), put
\[
  S''=\operatorname{Spec}(A'')
\]
and let \(X'',Z_i''\) be obtained from \(X,Z_i\) by the base change
\(S''\to S\). We obtain the diagram
\SGARefTarget{sga3:ix:lemma:4:4:diagram:01}
\[
\begin{tikzcd}[column sep=large,row sep=large]
  Z^i \arrow[d] & Z_n^{i''} \arrow[l]\arrow[d]
    & Z_n^{i'} \arrow[l]\arrow[d] \\
  X \arrow[d] & X'' \arrow[l]\arrow[d] & X' \arrow[l]\arrow[d] \\
  S & S'' \arrow[l] & S' \arrow[l].
\end{tikzcd}
\]

In what follows, we restrict to those \(A''\) that are localizations of
\(A\)-subalgebras \(A_1''\) of \(A'\), of finite type over \(A\), at
\(\mathfrak m'\cap A_1''\), where \(\mathfrak m'\) is the maximal ideal
of \(A'\). Thus the homomorphisms \(A\to A''\to A'\) are local. These
\(A''\) form an increasing filtered family of subalgebras whose union,
and hence whose direct limit, is \(A'\). Likewise,
\[
  B'=\varinjlim_{A''}\bigl(B\otimes_A A''\bigr).
\]
There therefore exist such an \(A''\) and an element
\(g''\in B''=B\otimes_A A''\) whose image in \(B'\) is \(g'\).

By \SGARefLink{sga3:ix:lemma:4:5}{Lemma~4.5 below}, the property that
\((Z_{is})_{i\in I}\) be schematically dense in \(X_s\) is preserved by
every base change \(S''\to S\).%
\footnote{Editor's note: the preceding sentence has been added, and the
original has been expanded in what follows.}
Consequently, because each \(S''\) is locally Noetherian, replacing
\(S\) by a suitable \(S''\) and \(X,Z_i\) by \(X'',Z_i''\) preserves
the hypotheses of \SGARefLink{sga3:ix:lemma:4:4}{Lemma~4.4}.

We may therefore suppose that \(g'\) comes from an element \(g\in B\).
After replacing \(X\) by the open subset \(X_g=U\), we may suppose that
\(U'=X'\). In the same way, we may suppose that \(f'\) comes from a
section \(f\) of \(\mathscr O_X\) over \(X\).

Let \(Y=V(f)\) be the subscheme of \(X\) defined by \(f\), and put
\[
  Y_i=Z_i\times_XY.
\]
Thus \(Y_i\) is the closed subscheme \(V(f_i)\) of \(Z_i\), where \(f_i\)
is the section of \(\mathscr O_{Z_i}\) induced by \(f\). Denoting by
\(Y',Y_i'\) the \(S'\)-schemes obtained by base change, we have
\SGARefTarget{sga3:ix:lemma:4:4:display:02}
\[
  Y'=V(f'),
  \qquad
  Y_i'=Z_i'\times_{X'}Y'=V(f_i'),
\]
and analogous relations hold for \(Y'',Y_i''\). The assertion that
\(f'\) vanishes on the \(Z_i'\) is equivalent to
\[
  Y_i'=Z_i'\qquad\text{for every }i.
\]
\footnote{Editor's note: the original added ``at least at every point
\(x'\) of \(X'\) above \(s'\),'' but this restriction seems unnecessary.}

Let \(\mathfrak m\) be the maximal ideal of \(A\). For every
\(n\geq0\), introduce the closed subscheme
\[
  S_n=\operatorname{Spec}(A/\mathfrak m^{n+1})
\]
of \(S\), and let \(X_n,Y_n,Z_n^i,Y_n^i\) be obtained from
\(X,Y,Z_i,Y_i\) by the base change \(S_n\to S\). More generally, for
every \(S\)-prescheme \(P\), write \(P_n=P\times_SS_n\).

For every \(n\) and \(i\in I\), consider the functor
\SGARefTarget{sga3:ix:lemma:4:4:display:03}
\[
  F_n^i
  =\prod_{Z_n^i/S_n}Y_n^i/Z_n^i:
  (\mathbf{Sch})_{/S_n}^{\circ}\longrightarrow\mathbf{Set},
\]
defined, for every \(P\) over \(S_n\), by
\SGARefTarget{sga3:ix:lemma:4:4:display:04}
\[
F_n^i(P)
=\Gamma\bigl((Y_n^i)_P/(Z_n^i)_P\bigr)
=
\begin{cases}
  \varnothing,&(Y_n^i)_P\ne(Z_n^i)_P,\\
  \{\operatorname{id}\},&(Y_n^i)_P=(Z_n^i)_P.
\end{cases}
\]
Since \(S_n\) is local Artinian and \(Z_n^i\) is flat, hence essentially
free, over \(S_n\), \SGARefLink{sga3:viii:theorem:6:4}{VIII, 6.4} shows
that each \(F_n^i\) is represented by a closed subprescheme \(T_n^i\)
of \(S_n\).

The completion \(\widehat A\) of the local ring \(A\) is Noetherian
because \(A\) is, and it is the projective limit of the rings \(A_n\).%
\footnote{Editor's note: the original has been expanded and corrected to
show that the projective limit of the rings
\(\mathscr O(T_n^i)\) defines a closed formal subscheme of
\(\operatorname{Spf}(\widehat A)\).}
Let \(K_n^i\) be the ideal of
\[
  A_n=A/\mathfrak m^{n+1}
\]
defining \(T_n^i\), and let \(K^i=\varprojlim_nK_n^i\); this is an ideal
of \(\widehat A\).

For fixed \(i\), integers \(m\geq n\), and any \(S_n\)-prescheme \(P\),
\[
  (Y_n^i)_P
  =Y_i\times_SS_n\times_{S_n}P
  =Y_i\times_SP
  =(Y_m^i)_P,
\]
and similarly \((Z_n^i)_P=(Z_m^i)_P\). It follows that \(F_n^i\) is the
restriction \(F_m^i\times_{S_m}S_n\) of \(F_m^i\) to
\((\mathbf{Sch})_{/S_n}\), and therefore
\[
  T_n^i=T_m^i\times_{S_m}S_n.
\]
We thus have a commutative diagram with exact rows:
\SGARefTarget{sga3:ix:lemma:4:4:diagram:02}
\[
\begin{tikzcd}[column sep=large,row sep=large]
  0 \arrow[r] & K_m^i \arrow[r]\arrow[d,two heads]
    & A_m \arrow[r]\arrow[d,two heads]
    & \mathscr O(T_m^i) \arrow[r]\arrow[d,two heads] & 0\\
  0 \arrow[r] & K_n^i \arrow[r]
    & A_n \arrow[r]
    & \mathscr O(T_n^i) \arrow[r] & 0
\end{tikzcd}
\]
All the vertical arrows are surjective. It follows, first, that the
projective system \((K_n^i)_{n\in\mathbb N}\) satisfies the
Mittag--Leffler condition, so that
\[
  \varprojlim_n\mathscr O(T_n^i)
  \simeq\widehat A/K^i
\]
as topological rings (cf. EGA \(0_{\mathrm{III}}\), \S13.2). Second, the
map \(K^i\to K_n^i\) is surjective (cf. \SGARefLink{sga3:ix:bib:bens}{[BEns]}, III, \S7.4,
Proposition~5). Hence
\[
  K_n^i\simeq
  \frac{K^i+\mathfrak m^{n+1}\widehat A}
       {\mathfrak m^{n+1}\widehat A},
  \qquad
  \mathscr O(T_n^i)\simeq
  (\widehat A/K^i)\otimes_A(A/\mathfrak m^{n+1}).
\]

In other words, \((T_n^i)_n\) is an inductive system of affine Artinian
schemes, and its inductive limit
\[
  T^i=\varinjlim_nT_n^i
\]
is a closed formal subscheme of
\[
  \widehat S=\operatorname{Spf}(\widehat A)
\]
(cf. EGA I, \S10), whose reduction modulo \(\mathfrak m^{n+1}\) is
\(T_n^i\).

Let \(T\) be the closed formal subscheme of \(\widehat S\) given by the
intersection of the \(T^i\), that is, defined by
\[
  K=\sum_iK^i.
\]
Since \(\widehat A\) is Noetherian, there is a finite subset \(J\) of
\(I\) such that
\[
  K=\sum_{i\in J}K^i,
  \qquad\text{equivalently}\qquad
  T=\bigcap_{i\in J}T^i.
\]
For every \(n\), therefore,
\[
  K_n=\sum_{i\in J}K_n^i,
\]
where \(K_n\) is the image of \(K\) in \(A_n\).

Recall that \(f_i\) is the image of
\(f\in B=\mathscr O(X)\) in \(\mathscr O(Z_i)\), and that \(f_i'\) is
its image in
\[
  \mathscr O(Z_i')
  =\mathscr O(Z_i)\otimes_AA'.
\]
The relation \(Y_i'=Z_i'\) is equivalent to \(f_i'=0\). Now
\(\mathscr O(Z_i)\otimes_AA'\) is the direct limit of the subalgebras
\(\mathscr O(Z_i)\otimes_AA''\), for \(A''\) as above. Consequently,
there is an \(A''\) for which \(Y_i''=Z_i''\). A priori \(A''\) depends
on \(i\), but, because \(J\) is finite, one \(A''\) can be chosen that
works for every \(i\in J\). Put
\[
  S''=\operatorname{Spec}(A''),
  \qquad
  S''_{(n)}=S''\times_SS_n.
\]
\footnote{Editor's note: the original has been expanded and corrected,
distinguishing \(S''_{(n)}\) from the subscheme
\(S_n''=S_n''(s'')\) introduced below.}

For every \(i\in J\), the equality
\[
  Y_i''\times_{S''}S''_{(n)}
  =(Y_n^i)_{S''_{(n)}}
  =Z_i''\times_{S''}S''_{(n)}
  =(Z_n^i)_{S''_{(n)}}
\]
shows, by the definition of \(T_n^i\), that \(S''_{(n)}\to S_n\)
factors through \(T_n^i\). It therefore factors through
\[
  T_n=\bigcap_{i\in J}T_n^i,
\]
and hence through \(T_n^i\) for every \(i\in I\). Let \(f''\) be the
image of \(f\) in \(B''=B\otimes_AA''\), and let \(f''_{(n)}\) be its
image in
\[
  B''_{(n)}=B''/\mathfrak m^{n+1}B''.
\]
The preceding factorization means that, for every \(n\), the section
\(f''_{(n)}\) vanishes on all the \(Z_i\times_SS''_{(n)}\).

Because \(A\to A''\to A'\) are local homomorphisms, the image \(s''\)
of \(s'\) in \(S''\) is the closed point of \(S''\), corresponding to
the maximal ideal \(\mathfrak m''\) of \(A''\), and its image in \(S\)
is \(s\). \emph{[The source prints \(B''\) here; since
\(s''\in S''=\operatorname{Spec}(A'')\), the intended ring is
\(A''\).]} Fix \(n\), and now let \(S_n''\) denote the \(n\)-th
infinitesimal neighborhood of \(s''\) in \(S''\):
\[
  S_n''=\operatorname{Spec}(A''/\mathfrak m''^{\,n+1}).
\]
Put
\[
  B_n''=B''/\mathfrak m''^{\,n+1}B'',
  \qquad
  Z_{n}^{i''}=Z_i\times_SS_n''.
\]
\emph{[The source appends
\(Z_n^{i''}=\operatorname{Spec}(B_n'')\); the displayed equality would
identify each \(Z_n^{i''}\) with \(X_n''\), so the base-change
definition is retained here.]}
Then the image \(f_n''\) of \(f''_{(n)}\) in \(B_n''\) vanishes on every
\(Z_n^{i''}\). By \SGARefLink{sga3:ix:lemma:4:5}{Lemma~4.5 below}, the
family of fibers
\[
  Z_0^{i'''}
  =Z_n^{i''}\times_{S_n''}\kappa(s'')
  =Z_{is}\times_{\kappa(s)}\kappa(s'')
\]
is schematically dense in
\[
  X_0'''
  =X_n''\times_{S_n''}\kappa(s'')
  =X_s\times_{\kappa(s)}\kappa(s'').
\]
Thus \(S_n'',X_n'',(Z_n^{i''})_{i\in I}\), and
\[
  S_0'''=\operatorname{Spec}\kappa(s'')
\]
satisfy the hypotheses of \SGARefLink{sga3:ix:lemma:4:2}{Lemma~4.2}.
It follows that \(f_n''=0\), that is,
\[
  f''\in\mathfrak m''^{\,n+1}B''
  \qquad\text{for every }n.
\]

Since \(B\) is Noetherian and \(B''=B\otimes_AA''\) is a localization
of a finite-type \(B\)-algebra, \(B''\) is Noetherian. Its local rings
are therefore separated for their usual topology. Hence \(f''\) is zero
at every point \(x''\) of \(X''\) for which
\(\mathfrak m''B''\) is contained in the maximal ideal of
\(\mathscr O_{X'',x''}\), that is, at every point of \(X''\) above
\(s''\). A fortiori, \(f'\) is zero at every point of \(X'\) above
\(s'\). \qed

\medskip
\noindent\SGARefTarget{sga3:ix:lemma:4:5:0}\textbf{Lemma 4.5.0.}%
\footnote{Editor's note: this lemma has been inserted here because it is
used in the proofs of \SGARefLink{sga3:ix:lemma:4:5}{4.5} and
\SGARefLink{sga3:ix:theorem:4:7}{4.7}.}
\textit{Let \(X\) be an \(S\)-prescheme, let \((Z_i)_{i\in I}\) be a
family of subpreschemes of \(X\), and let \(S'\to S\) be faithfully
flat. If \((Z_i')_{i\in I}\) is schematically dense in \(X'\), then
\((Z_i)_{i\in I}\) is schematically dense in \(X\).}

This follows from EGA IV\(_3\), 11.10.5(i) and 11.9.10(i). It remains to
give the proof of \SGARefLink{sga3:ix:lemma:4:5}{the following lemma}.

\medskip
\noindent\SGARefTarget{sga3:ix:lemma:4:5}\textbf{Lemma 4.5.}
\textit{Let \(X\) be a locally Noetherian prescheme}%
\footnote{Editor's note: this hypothesis is in fact superfluous; cf. EGA
IV\(_3\), 11.10.6 and 11.9.13.}
\textit{ over a field \(k\), let \((Z_i)_{i\in I}\) be a family of
subpreschemes of \(X\), let \(k'/k\) be a field extension, and let
\(X',Z_i'\) be obtained from \(X,Z_i\) by base change. Then
\((Z_i)_{i\in I}\) is schematically dense in \(X\) if and only if
\((Z_i')_{i\in I}\) is schematically dense in \(X'\).}

\emph{Proof.}
The ``if'' direction follows from
\SGARefLink{sga3:ix:lemma:4:5:0}{Lemma~4.5.0}; let us prove the
converse.%
\footnote{Editor's note: the original has been expanded in what follows.}
First, we may suppose \(X=\operatorname{Spec}(B)\) affine. For every
\(i\), let \((Z_{ij})_{j\in J_i}\) be an affine open covering of \(Z_i\).
Replacing \((Z_i)_{i\in I}\) by
\((Z_{ij})_{(i,j)\in J}\), where
\[
  J=\coprod_{i\in I}J_i,
\]
we may also suppose that every \(Z_i\) is affine.

Let
\[
  g\in B'=B\otimes_kk',
  \qquad
  t'\in B_g'
\]
be a section of \(\mathscr O_{X'}\) on the affine open
\(U'=X_g'\) whose restriction to each \(Z_i'\cap U'\) is zero. Thus its
image in every
\[
  \bigl(\mathscr O(Z_i)\otimes_kk'\bigr)_g
\]
is zero. There exists a finite-type \(k\)-subalgebra \(A\) of \(k'\)
such that
\[
  g\in B_A=B\otimes_kA,
  \qquad
  t'\in(B_A)_g.
\]
The map
\[
  \mathscr O(Z_i)\otimes_kA
  \longrightarrow
  \mathscr O(Z_i)\otimes_kk'
\]
is injective, and therefore so is
\[
  \bigl(\mathscr O(Z_i)\otimes_kA\bigr)_g
  \longrightarrow
  \bigl(\mathscr O(Z_i)\otimes_kk'\bigr)_g.
\]
The hypothesis consequently implies that the image of \(t'\) in every
\(\bigl(\mathscr O(Z_i)\otimes_kA\bigr)_g\) is zero. It remains to show
that \((Z_{iA})_{i\in I}\) is schematically dense in \(X_A\), which
follows from EGA IV\(_3\), 11.9.10(b).%
\footnote{Editor's note: the original continued: ``Using
\SGARefLink{sga3:ix:corollary:4:3}{4.3} for the
\(X_{A'}\), where \(A'\) is a quotient of \(A\) by a power of a maximal
ideal, and Hilbert's Nullstellensatz, one is reduced to the case where
\(k'\) is a finite extension of \(k\).'' This reduction could be
expanded and that argument pursued, since the case of a finite extension
\(k'/k\) is somewhat simpler than the more general case treated in EGA
IV\(_3\), 11.9.10(b).}
\qed

Let us record in passing \SGARefLink{sga3:ix:corollary:4:6}{the following result}, which will be used in a
later expos\'e.%
\footnote{Editor's note: this reference remains to be specified.}

\medskip
\noindent\SGARefTarget{sga3:ix:corollary:4:6}\textbf{Corollary 4.6.}
\textit{Let \(X\) be a flat \(S\)-prescheme and \(U\) an open subset of
\(X\). Suppose that \(U_s\) is schematically dense in \(X_s\) for every
\(s\in S\), and that either \(X\) is locally Noetherian or \(X\) is
locally of finite presentation over \(S\). Then \(U\) is schematically
dense in \(X\).}

\emph{Proof.}
For simplicity, suppose \(U\) is retrocompact in \(X\) (cf. EGA
IV\(_3\), 11.10.10 and 11.9.17 for the general case). The locally
Noetherian case is included only for reference; it is contained in
\SGARefLink{sga3:ix:corollary:4:3}{4.3}.

In the second case, we may plainly suppose \(S\) and \(X\) affine. Then
\(X\) and \(U\) are of finite presentation over
\(S=\operatorname{Spec}(A)\). The ring \(A\) is the direct limit of its
finite-type \(\mathbb Z\)-subalgebras \(A_i\). The ``patented
procedure'' already used (cf. EGA IV\(_3\), 8.8.2 and 8.10.5(iii))
shows that there are an index \(i\), an affine scheme \(X_i\) over
\[
  S_i=\operatorname{Spec}(A_i),
\]
and an open subset \(U_i\) of \(X_i\), from which \(X,U\) are obtained
by the base change \(S\to S_i\). Let \(E_i\) be the subset of \(S_i\)
consisting of those \(s\in S_i\) for which \((U_i)_s\) is schematically
dense in \((X_i)_s\).

By EGA IV\(_2\), 5.9.9 and 5.10.2, \(E_i\) is the set of points
\(s\in S_i\) for which \((U_i)_s\) contains the set
\[
  \operatorname{Ass}\mathscr O_{(X_i)_s}
\]
of points associated with the structure sheaf of \((X_i)_s\). By EGA
IV\(_3\), 11.9.17.1, this is a constructible condition; hence \(E_i\)
is constructible.%
\footnote{Editor's note: the original has been expanded in this paragraph
and the next.}

By \SGARefLink{sga3:ix:lemma:4:5}{4.5}, the inverse image of \(E_i\)
under \(S_j\to S_i\) is \(E_j\), and its inverse image under
\(S\to S_i\) is the set \(E\) of those \(s\in S\) for which \(U_s\)
is schematically dense in \(X_s\).%
\footnotemark[24]
By hypothesis, \(E=S\), which is also the inverse image of \(S_i\) under
every \(S\to S_i\). EGA IV\(_3\), 8.3.11 therefore implies that there
is a \(j\geq i\) such that \(E_j=S_j\); that is, \((U_j)_s\) is
schematically dense in \((X_j)_s\) for every \(s\in S_j\).

Apply \SGARefLink{sga3:ix:lemma:4:4}{4.4} to
\((S_j,X_j,U_j)\) and the base change \(S\to S_j\). It follows that
\(U\) is schematically dense in \(X\). Notice that here
\SGARefLink{sga3:ix:lemma:4:4}{4.4} is used only
for a family with a finite index set (in fact, a singleton); in this
case its proof simplifies considerably. \qed

\medskip
\noindent\SGARefTarget{sga3:ix:theorem:4:7}\textbf{Theorem 4.7.}
\textit{Let \(S\) be a prescheme, and let \(G\) be a group of
multiplicative type and of finite type over \(S\). For every integer
\(n>0\), let \({}_nG\) be the kernel of \(n\cdot\operatorname{id}_G\).
Then the family of subpreschemes \(({}_nG)_{n>0}\) is schematically
dense in \(G\) in the sense of
\SGARefLink{sga3:ix:definition:4:1}{Definition~4.1}.}

\emph{Proof.}
We distinguish two cases.

\smallskip
\SGARefTarget{sga3:ix:theorem:4:7:proof:item:a}\noindent\textit{(a) \(S\) locally Noetherian.}
By \SGARefLink{sga3:ix:corollary:4:3}{4.3}, we are reduced to the case
in which \(S\) is the spectrum of a field \(k\). By
\SGARefLink{sga3:ix:lemma:4:5:0}{4.5.0}, we may suppose that \(k\) is
algebraically closed and that \(G\) is diagonalizable, say
\[
  G=D_k(M),
\]
where \(M\) is a commutative group of finite type. Then
\[
  M\simeq\Gamma\times\mathbb Z^r
\]
with \(\Gamma\) finite, and consequently
\[
  G\simeq G_1\times G_2,
  \qquad
  G_1=D(\Gamma),
  \qquad
  G_2=\mathbf G_m^r.
\]
For \(n\) sufficiently divisible---namely, for \(n\) a multiple of the
order of \(\Gamma\)---one has
\[
  {}_nG=G_1\times{}_nG_2,
  \qquad\text{since}\qquad
  {}_nG_1=G_1.
\]
Applying \SGARefLink{sga3:ix:corollary:4:3}{4.3} again to the projection
\(G\to G_1\), we are reduced to
\(G_2\), hence to the case \(G=\mathbf G_m^r\). Since \(G\) is then
reduced, schematic density of \(({}_nG)_{n>0}\) is equivalent to
density of \(\bigcup_n{}_nG\) in the ordinary topology. As
\[
  {}_nG=({}_n\mathbf G_m)^r,
\]
we are reduced to \(G=\mathbf G_m\), an irreducible group of dimension
one. The assertion now follows because
\(\bigcup_n{}_nG\), the set of roots of unity in \(k\), is infinite.

\smallskip
\SGARefTarget{sga3:ix:theorem:4:7:proof:item:b}\noindent\textit{(b) The general case.}
For every \(s\in S\), there are an open neighborhood \(U\) of \(s\)
and a faithfully flat quasi-compact morphism \(S'\to U\) such that
\(G'=G_{S'}\) is diagonalizable, say \(G'=D_{S'}(M)\). By
\SGARefLink{sga3:ix:lemma:4:5:0}{4.5.0} we may
reduce to \(G'\), and therefore suppose that \(G\) itself is
diagonalizable. It is obtained by base change
\[
  S\longrightarrow\operatorname{Spec}(\mathbb Z)
\]
from the absolute group \(H=D_{\mathbb Z}(M)\). Part~\SGARefLink{sga3:ix:theorem:4:7:proof:item:a}{(a)} shows that,
for every \(s\in\operatorname{Spec}(\mathbb Z)\), the family
\(({}_nH_s)_{n>0}\) is schematically dense in \(H_s\). It remains only
to apply \SGARefLink{sga3:ix:lemma:4:4}{4.4}.%
\footnote{Editor's note: \SGARefLink{sga3:ix:theorem:4:7}{this result} will be needed in X, \SGARefLink{sga3:x:lemma:4-3}{4.3} for a
non-locally-Noetherian base, namely
\(S=\operatorname{Spec}(\widehat A\otimes_A\widehat A)\), where
\(\widehat A\) is the completion of the Noetherian local ring \(A\).}
\qed

\medskip
\noindent\SGARefTarget{sga3:ix:corollary:4:8}\textbf{Corollary 4.8.}
\begin{enumerate}
  \item[(a)]\SGARefTarget{sga3:ix:corollary:4:8:item:a}
  \textit{Let \(u,v:H\to G\) be two homomorphisms of \(S\)-preschemes
  in groups, with \(H\) of multiplicative type}%
  \footnote{Editor's note: ``diagonalizable'' has been replaced by ``of
  multiplicative type.''}
  \textit{ and of finite type. Suppose that, for every integer \(n>0\),
  the restrictions of \(u\) and \(v\) to \({}_nH\) are equal. Then
  \(u=v\).}

  \item[(b)]\SGARefTarget{sga3:ix:corollary:4:8:item:b}
  \textit{Let \(H_1,H_2\) be two subgroups of multiplicative type and
  of finite type of a prescheme in groups \(G\). If
  \({}_nH_1={}_nH_2\) for every integer \(n>0\), then \(H_1=H_2\).}
\end{enumerate}

\emph{Proof.}
Part~\SGARefLink{sga3:ix:corollary:4:8:item:a}{(a)} follows from part~\SGARefLink{sga3:ix:corollary:4:8:item:b}{(b)} by considering the subgroups \(H_1,H_2\)
of \(H\times G\) that are the graphs of \(u,v\). For part~\SGARefLink{sga3:ix:corollary:4:8:item:b}{(b)}, put
\[
  H=H_1\cap H_2.
\]
This is a subgroup prescheme of \(H_i\), for \(i=1,2\), and we must show
that it equals \(H_i\). The hypothesis says precisely that \(H\)
contains all the \({}_nH_i\). We are therefore reduced to \SGARefLink{sga3:ix:corollary:4:9}{the next
corollary}. \qed

\medskip
\noindent\SGARefTarget{sga3:ix:corollary:4:9}\textbf{Corollary 4.9.}
\textit{Let \(G\) be an \(S\)-group of multiplicative type and of finite
type, and let \(H\) be a subgroup prescheme containing \({}_nG\) for
every \(n>0\). Then \(H=G\).}

\emph{Proof.}
By \SGARefLink{sga3:ix:theorem:4:7}{4.7}, it is enough to prove that
\(H\) is closed, or equivalently that it equals \(G\) set-theoretically.
This reduces us to the case in which \(S\) is the spectrum of a field.
Every subgroup prescheme of \(G\) is then closed, by VI\(_B\), \SGARefLink{cor:sga3-VIB-1.4.2}{1.4.2}.
\qed

\medskip
\noindent\SGARefTarget{sga3:ix:remark:4:10}\textbf{Remark 4.10.}
Under the hypotheses of \SGARefLink{sga3:ix:theorem:4:7}{4.7}, let
\(m>0\) be an integer with the following properties: for every
\(s\in S\), \(m\) is not a power of the characteristic of \(\kappa(s)\);
and if \(G_s\) has type \(M\), every prime divisor of the torsion
subgroup of \(M\) divides \(m\). The second condition is automatically
satisfied if \(G\) is a torus. The proof shows that in
\SGARefLink{sga3:ix:theorem:4:7}{Theorem~4.7} and
\SGARefLink{sga3:ix:corollary:4:8}{Corollary~4.8} and
\SGARefLink{sga3:ix:corollary:4:9}{Corollary~4.9} one may restrict to
the subgroups
\[
  {}_{m^r}G,\qquad r>0.
\]

Moreover, if \(G\) is smooth over \(S\), then for every \(s\in S\) there
are an open neighborhood \(U\) of \(s\) and an integer \(m>0\), prime
to every residue characteristic on \(U\), that satisfy the preceding
conditions for \(G_U\). For instance, take \(m\) to be the order of the
torsion subgroup of the type of \(G\) at \(s\) (cf.
\SGARefLink{sga3:ix:remark:1:4:1:item:a}{1.4.1(a)} and
\SGARefLink{sga3:ix:proposition:2:1:item:e:01}{2.1(e)}). After restriction to
\(U\), the groups \({}_{m^r}G\) are then
finite and \'etale over \(S\), and their family is schematically dense.
In certain applications---in particular those involving
\SGARefLink{sga3:ix:theorem:3:2}{Theorem~3.2} and
\SGARefLink{sga3:ix:theorem:3:2:bis}{Theorem~3.2~bis} (cf. XI, 6),
though not those involving \SGARefLink{sga3:ix:theorem:3:6}{Theorem~3.6}
and \SGARefLink{sga3:ix:theorem:3:6:bis}{Theorem~3.6~bis}---this
observation allows one to avoid
\SGARefLink{sga3:ix:theorem:3:1}{Theorem~3.1} and its use of Hochschild
cohomology.

\SGARefTarget{sga3:ix:section:5}\section*{5. Central Homomorphisms of Groups of Multiplicative Type}
\addcontentsline{toc}{section}{5. Central Homomorphisms of Groups of Multiplicative Type}

\medskip
\noindent\SGARefTarget{sga3:ix:lemma:5:0}\textbf{Lemma 5.0.}%
\footnote{Editor's note: this lemma has been made explicit because it is
used several times below; it was implicit in the original.}
\textit{Let \((A,\mathfrak m)\) be a Noetherian local ring, put
\(S=\operatorname{Spec}(A)\), and let \(H\) be a finite scheme over
\(S\), say \(H=\operatorname{Spec}(B)\), where \(B\) is finite over
\(A\). Let \(K\) be a subscheme of \(H\) such that reduction modulo
\(\mathfrak m^{n+1}\) gives \(K_n=H_n\) for every \(n\). Then \(K=H\).}

\emph{Proof.}
Let \(s\) be the closed point of \(S\). First, \(K\) is a closed
subscheme of \(H\). Indeed, a priori it is a closed subscheme of an open
subset \(U\) of \(H\). But \(K\), and hence \(U\), contains the fiber
\[
  K_s=H_s.
\]
Since \(H\to S\) is finite, and therefore closed, the complement of
\(U\) is empty; thus \(U=H\). Hence \(K\) is defined by an ideal \(I\)
of \(B\). The hypothesis implies
\[
  I\subset\mathfrak m^nB
  \qquad\text{for every }n.
\]
Since \(B\) is a finite \(A\)-module, it is separated for the
\(\mathfrak m\)-adic topology. Therefore \(I=0\). \qed

\medskip
\noindent\SGARefTarget{sga3:ix:theorem:5:1}\textbf{Theorem 5.1.}
\textit{Let \(u,v:H\to G\) be two homomorphisms of \(S\)-preschemes in
groups, where \(H\) is of multiplicative type and of finite type.
Suppose either that \(S\) is locally Noetherian or that \(G\) is of
finite presentation over \(S\). Let \(s\in S\) satisfy \(u_s=v_s\), and
suppose that \(u_s\) is central. Then there is an open neighborhood
\(U\) of \(s\) such that \(u_U=v_U\).}

\emph{Proof.}
We distinguish the two cases.

\smallskip
\SGARefTarget{sga3:ix:theorem:5:1:proof:item:a}\noindent\textit{(a) \(S\) locally Noetherian.}%
\footnote{Editor's note: the original has been expanded in what follows.}
Let
\[
  K=\operatorname{Ker}(u,v)
\]
be the inverse image of the diagonal of \(G\times_SG\) under
\((u,v):H\to G\times_SG\). It is a subgroup prescheme of \(H\). We seek
an open \(U\) for which \(K_U=H_U\). Because \(S\) is locally
Noetherian and \(H\) is of finite type over \(S\), the prescheme \(H\)
is locally Noetherian. Hence the immersion \(K\hookrightarrow H\) is of
finite type (cf. EGA I, 6.3.5), so \(K\) is of finite type, and therefore
of finite presentation, over \(S\). By EGA IV\(_3\), 8.8.2.4, it is
enough to prove that
\[
  K_{S_0}=H_{S_0},
  \qquad
  S_0=\operatorname{Spec}(\mathscr O_{S,s}).
\]
We may thus suppose that \(S\) is local with closed point \(s\).
Replacing the Noetherian local ring \(A\) by its completion
\(\widehat A\), if desired, gives a faithfully flat quasi-compact base
change \(\widehat S\to S\), and allows us to suppose that \(A\) is
complete.%
\footnote{Editor's note: this is because \(K=H\) if
\(K_{\widehat S}=H_{\widehat S}\); cf. SGA~1, VIII, \SGARefLink{sga3:viii:corollary:5:4}{5.4}. The argument
below does not use completeness of \(A\).}

By \SGARefLink{sga3:ix:corollary:3:4}{3.4}, one has \(u_n=v_n\) for every
\(n\), where the subscript denotes reduction modulo
\(\mathfrak m^{n+1}\), with \(\mathfrak m\) the maximal ideal of \(A\).
For every \(m>0\), let
\[
  {}_mu,{}_mv:{}_mH\longrightarrow G
\]
be the homomorphisms induced by \(u,v\). Then
\[
  ({}_mu)_n=({}_mv)_n
  \qquad\text{for every }n.
\]
The group \({}_mH\) is finite over \(S\), by \SGARefLink{sga3:ix:proposition:2:2}{2.2}, so
\SGARefLink{sga3:ix:lemma:5:0}{5.0} gives \({}_mu={}_mv\). This holds
for every \(m\), and therefore
\SGARefLink{sga3:ix:theorem:4:7}{4.7} gives \(u=v\).

\smallskip
\SGARefTarget{sga3:ix:theorem:5:1:proof:item:b}\noindent\textit{(b) \(G\) of finite presentation.}%
\footnote{Editor's note: the original has been expanded in what follows.}
The group \(H\) is also of finite presentation over \(S\). By EGA
IV\(_3\), 8.8.2, we may again suppose that \(S\) is local with closed
point \(s\), and it suffices to prove \(u=v\). Let \(f:S'\to S\) be a
faithfully flat quasi-compact morphism. Write
\[
  f_H:H'\to H,
  \qquad
  f_G:G'\to G,
  \qquad
  u',v':H'\to G'
\]
for the resulting morphisms. If \(u'=v'\), then
\[
  u\circ f_H=f_G\circ u'=f_G\circ v'=v\circ f_H,
\]
and \(u=v\) because \(f_H\) is an epimorphism. After a faithfully flat
quasi-compact extension of the base, we may therefore suppose that
\(H\) is diagonalizable:
\[
  H=D_S(M)
\]
for a commutative group \(M\) of finite type.

As in the proof of \SGARefLink{sga3:ix:corollary:4:6}{4.6}, introduce
the increasing filtered family of \(\mathbb Z\)-subalgebras \(A_i\) of
\(A=\mathscr O_{S,s}\) that are of finite type, and put
\[
  S_i=\operatorname{Spec}(A_i).
\]
\footnotemark[30]
For every \(i\), the group \(H=D_S(M)\) comes from
\[
  H_i=D_{S_i}(M).
\]
Since \(G\) is of finite presentation over \(S\), EGA IV\(_3\), 8.8.2
(see also VI\(_B\), 10.2 and 10.3) gives an index \(i\), an
\(S_i\)-group prescheme \(G_i\) of finite presentation, and
homomorphisms
\[
  u_i,v_i:H_i\longrightarrow G_i
\]
whose base changes are \(u,v\). Let \(s_i\) be the image of \(s\) in
\(S_i\), and define
\[
  \rho_i:H_i\times_{S_i}G_i\longrightarrow G_i,
  \qquad
  \rho_i(h,g)=u_i(h)g\,u_i(h)^{-1}.
\]
Because \(u_s\) is central,
\[
  \rho_s=\rho_i\times_{\kappa(s_i)}\kappa(s)
\]
is the second projection. The extension
\(\kappa(s_i)\to\kappa(s)\) is faithfully flat, so \(\rho_i\) is also
the second projection over \(s_i\); hence \((u_i)_{s_i}\) is central.
Likewise \(u_s=v_s\) implies
\[
  (u_i)_{s_i}=(v_i)_{s_i}.
\]
Part~\SGARefLink{sga3:ix:theorem:5:1:proof:item:a}{(a)}, applied over \(S_i\), now gives the conclusion. \qed

\medskip
\noindent\SGARefTarget{sga3:ix:corollary:5:2}\textbf{Corollary 5.2.}
\textit{Let \(u:H\to G\) be a homomorphism of \(S\)-preschemes in
groups, where \(H\) is of multiplicative type and of finite type, and
suppose either that \(S\) is locally Noetherian or that \(G\) is of
finite presentation over \(S\). If \(s\in S\) and \(u_s\) is the unit
homomorphism, then there is an open neighborhood \(U\) of \(s\) on
which \(u_U\) is the unit homomorphism.}

\medskip
\noindent\SGARefTarget{sga3:ix:corollary:5:3}\textbf{Corollary 5.3.}
\textit{Let \(u,H,G\) be as in
\SGARefLink{sga3:ix:corollary:5:2}{5.2}, and suppose in addition that \(G\)
is separated over \(S\). Let \(U\) be the set of points \(s\in S\) for
which \(u_s\) is the unit homomorphism. Then \(U\) is open and closed
in \(S\), and \(u_U:H_U\to G_U\) is the unit homomorphism.}

\emph{Proof.}
Only closedness remains to be proved. Put \(K=\operatorname{Ker}(u)\).
Since \(G\) is separated over \(S\), \(K\) is a closed subprescheme of
\(H\), and \(U\) is the set of points \(s\in S\) for which \(K_s=H_s\).
It follows readily that \(U\) is closed, for example by
\SGARefLink{sga3:viii:theorem:6:4}{VIII, 6.4}, since \(H\) is
essentially free over \(S\) by VIII, \SGARefLink{sga3:viii:proposition:6:3}{6.3}. Alternatively, we may suppose
that \(S\) is reduced and that \(U\) is dense in \(S\), hence
schematically dense. Then \(H_U\) is schematically dense in \(H\)
because \(H\) is flat over \(S\),%
\footnote{Editor's note: and of finite presentation over \(S\); cf. EGA
IV\(_3\), 11.10.5(ii) and 11.9.10(ii).}
and \(u\) and the unit homomorphism, being equal on \(H_U\), are equal
on \(H\). \qed

\medskip
\noindent\SGARefTarget{sga3:ix:corollary:5:4}\textbf{Corollary 5.4.}
\textit{Let \(S\) be a prescheme; let \(H,G\) be \(S\)-groups, with
\(H\) of multiplicative type and of finite type and \(G\) separated
and of finite presentation over \(S\); and let
\(\pi:S'\to S\) be a}%
\footnote{Editor's note: ``covering morphism for the faithfully flat
quasi-compact topology'' has been replaced by ``universal effective
epimorphism,'' so that the result applies to the morphism \(K\to S\) in
\SGARefLink{sga3:ix:corollary:5:5}{5.5}; the beginning of the proof has
been modified accordingly.}
\textit{ universal effective epimorphism---for example, a faithfully
flat quasi-compact morphism or a morphism admitting a section (cf. IV,
\SGARefLink{sga3:iv:proposition:1:12}{1.12})---whose fibers are geometrically connected.}%
\footnote{Editor's note: this holds, for example, when
\(S=\operatorname{Spec}(k)\) and \(S'=\operatorname{Spec}(k')\), where
\(k\) is a field and \(k'/k\) is radicial; cf. X, \SGARefLink{sga3:x:proposition:1:4}{Proposition~1.4}.}

\textit{Let \(u':H'\to G'\) be a central homomorphism of the
\(S'\)-groups obtained from \(H,G\) by base change. Then there is a
unique homomorphism \(u:H\to G\) of \(S\)-groups such that
\(u\times_SS'=u'\). If \(S'/S\) admits a section \(g\), then \(u\) is
the morphism obtained from \(u'\) by the base change
\(g:S\to S'\).}

\emph{Proof.}
Because \(\pi\) is a universal effective epimorphism, so is
\(H'\to H\); uniqueness follows as at the beginning of the proof of
\SGARefLink{sga3:ix:theorem:5:1:proof:item:b}{5.1(b)}. If \(\pi\) has a section
\(g\), then
\[
  u'=\pi^*u
  \quad\Longrightarrow\quad
  u=g^*\pi^*u=g^*u'.
\]

For existence, IV, \SGARefLink{sga3:iv:proposition:2:3}{2.3} reduces us to showing that the two homomorphisms
\[
  u_1'',u_2'':H''\longrightarrow G''
\]
of \(S''\)-groups obtained from \(u'\) by the two base changes
\[
  \operatorname{pr}_1,\operatorname{pr}_2:
  S''=S'\times_SS'\longrightarrow S'
\]
are equal. Their pullbacks along the diagonal \(S'\to S''\) are equal,
since both are \(u'\).%
\footnote{Editor's note: the preceding sentence has been added.}
Because \(u_1'',u_2''\) are central, we may apply
\SGARefLink{sga3:ix:corollary:5:3}{5.3} to
\[
  u_1''(u_2'')^{-1}.
\]
There is therefore an open-and-closed subset \(U\) of \(S''\),
containing the diagonal, over which \(u_1''\) and \(u_2''\) coincide.
The fibers of \(S'/S\) are geometrically connected, hence so are those
of \(S''/S\), and in particular they are connected. Since \(U\)
contains the diagonals of these fibers, it contains every fiber and
therefore equals \(S''\). This proves existence. \qed

\medskip
\noindent\SGARefTarget{sga3:ix:corollary:5:5}\textbf{Corollary 5.5.}
\textit{Let \(S\) be a prescheme, let \(K\) be an \(S\)-group prescheme
locally of finite type}%
\footnote{Editor's note: the hypothesis ``locally of finite type'' has
been added to agree with VI\(_A\), \SGARefLink{sga3:ix:corollary:2:4}{2.4}. In fact, over a field every
connected group scheme is geometrically connected.}
\textit{ and with connected fibers, and let \(H\) be an invariant
subgroup of \(K\), of multiplicative type and of finite type. Then
\(H\) is a central subgroup of \(K\).}

\emph{Proof.}%
\footnote{Editor's note: the original has been expanded in what follows.}
First, \(\pi:K\to S\) has geometrically connected fibers, because for a
group scheme locally of finite type over a field, connectedness implies
geometric connectedness (cf. VI\(_A\), \SGARefLink{sga3:ix:corollary:2:4}{2.4}). Apply
\SGARefLink{sga3:ix:corollary:5:4}{5.4} with \(G=H\) and \(S'=K\) to
the homomorphism of \(K\)-groups
\[
  u':H_K\longrightarrow H_K,
  \qquad
  (h,k)\longmapsto(khk^{-1},k).
\]
It is central because \(H\) is commutative. Its pullback along the unit
section \(\epsilon:S\to K\) is the identity homomorphism of \(H\).
\emph{[The source prints ``of \(K\)'' here; the displayed endomorphism
of \(H_K\) shows that \(H\) is intended.]}
\SGARefLink{sga3:ix:corollary:5:4}{Corollary~5.4} therefore implies
that \(u'\) itself is the identity on
\(H_K\). Thus \(H\) is central in \(K\). \qed

Let us state the variants of the preceding results for central subgroups
of multiplicative type. Arguing as above, and using
\SGARefLink{sga3:ix:theorem:3:2:bis}{3.2~bis}, gives:

\medskip
\noindent\SGARefTarget{sga3:ix:theorem:5:1:bis}\textbf{Theorem 5.1 bis.}
\textit{Let \(G\) be an \(S\)-group prescheme, and let \(H_1,H_2\) be
two subgroup preschemes of multiplicative type and of finite type, with
\((H_1)_s\) central. Suppose either that \(S\) is locally Noetherian or
that \(G\) is of finite presentation over \(S\). For every \(s\in S\)
such that \((H_1)_s=(H_2)_s\), there is an open neighborhood \(U\) of
\(s\) for which \(H_1|_U=H_2|_U\).}

\medskip
\noindent\SGARefTarget{sga3:ix:corollary:5:3:bis}\textbf{Corollary 5.3 bis.}
\textit{Under the preceding hypotheses, suppose moreover that \(G\) is
separated over \(S\). Let \(U\) be the set of points \(s\in S\) for
which \((H_1)_s=(H_2)_s\). Then \(U\) is open and closed in \(S\), and
\(H_1|_U=H_2|_U\).}

\medskip
\noindent\SGARefTarget{sga3:ix:corollary:5:4:bis}\textbf{Corollary 5.4 bis.}
\textit{Let \(S\) be a prescheme, let \(G\) be a group prescheme
separated and of finite presentation over \(S\), let \(S'\to S\) be a
covering morphism for the faithfully flat quasi-compact topology with
geometrically connected fibers, and let \(H'\) be a subgroup of
multiplicative type and of finite type of \(G'=G\times_SS'\).
Then there is a unique subgroup \(H\) of \(G\) such that
\(H'=H\times_SS'\). The subgroup \(H\) is of multiplicative type and
of finite type over \(S\), by \SGARefLink{sga3:ix:definition:1:1}{1.1} and \SGARefLink{sga3:ix:proposition:2:1:item:b:01}{2.1(b)}. If \(S'\to S\) admits a
section \(g\), then \(H=g^*H'\).}

\medskip
\noindent\SGARefTarget{sga3:ix:theorem:5:6}\textbf{Theorem 5.6.}
\textit{Let \(S\) be a prescheme and let \(u:H\to G\) be a homomorphism
of \(S\)-preschemes in groups, where \(H\) is of multiplicative type and
of finite type and \(G\) is of finite presentation over \(S\).}
\begin{enumerate}
  \item[(a)]\SGARefTarget{sga3:ix:theorem:5:6:item:a}
  \textit{Suppose that \(G\) has connected fibers. If \(s\in S\) and
  \(u_s:H_s\to G_s\) is central, then there is an open neighborhood
  \(U\) of \(s\) such that the homomorphism \(H|_U\to G|_U\) induced by
  \(u\) is central.}

  \item[(b)]\SGARefTarget{sga3:ix:theorem:5:6:item:b}
  \textit{If \(u_s:H_s\to G_s\) is central for every \(s\in S\), then
  \(u\) is central.}
\end{enumerate}

\emph{Proof.}
We first prove \SGARefLink{sga3:ix:theorem:5:6:item:a}{(a)}.%
\footnote{Editor's note: the original has been expanded in the proof of
\SGARefLink{sga3:ix:theorem:5:6:item:a}{(a)}.}
Arguing exactly as in the proof of
\SGARefLink{sga3:ix:theorem:5:1:proof:item:b}{5.1(b)}, we may suppose first that
\(S\) is local with closed point \(s\), then that
\(H=D_S(M)\) is diagonalizable, and finally that there exist a
finite-type \(\mathbb Z\)-subalgebra \(A_i\) of
\(A=\mathscr O_{S,s}\) and a homomorphism
\[
  u_i:D_{S_i}(M)\longrightarrow G_i
\]
whose base change is \(u\) and for which \((u_i)_{s_i}\) is central,
where \(s_i\) is the image of \(s\) in \(S_i\). We are thereby reduced
to the case in which \(S\) is Noetherian local with closed point \(s\),
and we must prove that \(u\) is central.

Let \(K=\operatorname{Ker}(v,w)\), where
\[
  v,w:H\times_SG\rightrightarrows G
\]
are defined by
\SGARefTarget{sga3:ix:theorem:5:6:display:01}
\[
  v(h,g)=g,
  \qquad
  w(h,g)=\operatorname{int}(u(h))\cdot g
        =u(h)g\,u(h)^{-1}.
\]
Then \(K\) is a subgroup over \(G\) of
\[
  H_G=H\times_SG,
\]
and we must prove that \(K=H_G\). By
\SGARefLink{sga3:ix:corollary:4:9}{4.9}, it is enough to show that \(K\)
contains
\[
  {}_m(H_G)=({}_mH)\times_SG
\]
for every \(m>0\). We are thus reduced to the case \(H={}_mH\), so that
\(H\) is finite over \(S\).

Let \(e\) be the unit element of the fiber \(G_s\), and put
\[
  S^0=\operatorname{Spec}(\mathscr O_{G,e}),
  \qquad
  S_n^0=S^0\times_SS_n,
  \qquad
  S_n=\operatorname{Spec}(A/\mathfrak m^{n+1}).
\]
The scheme \(S^0\) is Noetherian local because \(G\) is of finite
presentation over the Noetherian scheme \(S\). The base change
\[
  K_{S^0}=K\times_GS^0
\]
is a subscheme of
\[
  H_{S^0}=H\times_SS^0,
\]
and \SGARefLink{sga3:ix:corollary:3:9}{3.9} gives
\[
  K_{S_n^0}=H_{S_n^0}
  \qquad\text{for every }n.
\]
\footnote{Editor's note: the original has been expanded and corrected,
replacing ``the unit section of \((G_H)_n\) over \(H_n\)'' by
\(H_{S_n^0}\).}
Since \(H_{S^0}\) is finite over \(S^0\),
\SGARefLink{sga3:ix:lemma:5:0}{5.0}, applied to the Noetherian local
ring \(\mathscr O_{G,e}\), gives
\[
  K_{S^0}=H_{S^0}.
\]
\emph{[The source prints the undefined local ring \(B\) here; the
Noetherian local ring \(\mathscr O_{G,e}\) just introduced is intended.]}

On the other hand, \(H\times_SG\) is Noetherian, being of finite
presentation over the Noetherian scheme \(S\). Thus the immersion
\[
  K\hookrightarrow H\times_SG
\]
is of finite type (cf. EGA I, 6.3.5), and \(K\) is of finite
presentation over \(G\). The equality
\[
  K_{S^0}=H_G\times_GS^0
\]
therefore implies, by EGA IV\(_3\), 8.8.2.4, that there is an open
neighborhood \(W\) of \(e\) in \(G\) such that
\[
  K\times_GW
  =H_G\times_GW
  =H\times_SW.
\]
Hence \(K\) contains the open neighborhood
\[
  V=H\times_SW
\]
of the unit section of \(G_H=G\times_SH\) over \(H\).

For every \(t\in H\), the fiber \(G_t\) is an algebraic group over
\(\kappa(t)\), hence Cohen--Macaulay by VI\(_A\), 1.1.1 and therefore
has no embedded components. It is also connected, hence irreducible by
VI\(_A\), 2.4, so its generic point is its unique associated point.
EGA IV\(_2\), 3.1.8 therefore shows that \(V_t\) is schematically dense
in \(G_t\). By \SGARefLink{sga3:ix:corollary:4:3}{4.3}, \(V\) is
schematically dense in \(G_H\).

Moreover, since \(K\) is an \(H\)-subgroup of \(G_H\), it induces on
each fiber \(G_h\) a subgroup \(K_h\). This subgroup contains an open
neighborhood of the unit element; because \(G_h\) is connected, it
follows that \(K_h=G_h\). Thus \(K=G_H\) set-theoretically. Finally,
\(K\) is a closed subprescheme of \(G_H\) containing the schematically
dense subprescheme \(V\), and hence \(K=G_H\). This proves \SGARefLink{sga3:ix:theorem:5:6:item:a}{(a)}.

Part~\SGARefLink{sga3:ix:theorem:5:6:item:b}{(b)} follows directly from
\SGARefLink{sga3:ix:corollary:3:9}{3.9}: to verify that a subprescheme
\(K\) of \(P=G\times_SH\) equals \(P\), it is enough to verify, after
every base change \(S'\to S\) in which \(S'\) is the spectrum of an
Artinian quotient of a local ring of \(S\), that \(K'=P'\). \qed

